\section{Quadratische Matrizen}

Eine Matrix $A$ heisst \emph{quadratisch}, wenn $m = n$. Die Elemente $a_{11}, a_{22}, \ldots, a_{nn}$ bilden die \emph{Hauptdiagonale}.

\subsection{Spezielle Matrizen}
\begin{itemize}
    \item \textbf{Diagonalmatrix}: Alle Elemente ausserhalb der Hauptdiagonalen $= 0$.
    \item \textbf{Einheitsmatrix} $E$: Diagonalmatrix mit allen Diagonalelementen $= 1$.
    \item \textbf{Obere Dreiecksmatrix}: $i > j \Rightarrow u_{ij} = 0$ (unterhalb $= 0$).
    \item \textbf{Untere Dreiecksmatrix}: $i < j \Rightarrow l_{ij} = 0$ (oberhalb $= 0$).
    \item \textbf{Symmetrische Matrix}: $S^T = S$.
\end{itemize}

\textbf{Einheitsmatrix:} $AE = EA = A$ \quad ($E$ spielt die Rolle der $1$).

\subsection{Potenzen}
$A^0 = E$, \quad $A^k = \underbrace{A \cdot A \cdots A}_{k}$ \quad ($k \in \mathbb{N}_{\ge 1}$)

Für Diagonalmatrizen: $D^k = \operatorname{diag}(a_{11}^k, \ldots, a_{nn}^k)$.

Im Allgemeinen: $(AB)^2 = ABAB \neq AABB = A^2B^2$ (weil Multiplikation nicht kommutativ).

\subsection{Inverse Matrizen}

Die \emph{Inverse} einer quadratischen Matrix $A$ ist $A^{-1}$ mit:
\[
    A \cdot A^{-1} = A^{-1} \cdot A = E
\]

\textbf{Inverse einer $2 \times 2$-Matrix:}
\[
    \begin{pmatrix} a & b \\ c & d \end{pmatrix}^{-1}
    = \frac{1}{ad - bc} \begin{pmatrix} d & -b \\ -c & a \end{pmatrix}
\]
Existiert nur, wenn $ad - bc \neq 0$.

\subsubsection{Invertierbar / Singulär}
\emph{Invertierbar (regulär)}: Matrix hat eine Inverse. \\
\emph{Singulär}: Matrix hat keine Inverse.

\subsubsection{Eigenschaften invertierbarer Matrizen}
\begin{enumerate}
    \item $A^{-1}$ ist eindeutig bestimmt.
    \item $(A^{-1})^{-1} = A$
    \item $(A \cdot B)^{-1} = B^{-1} \cdot A^{-1}$ \quad (Reihenfolge umkehren!)
    \item $(A^T)^{-1} = (A^{-1})^T$
\end{enumerate}

\subsubsection{Zusammenhang mit LGS}
Für $A \cdot \vec{x} = \vec{c}$ mit $n \times n$-Matrix $A$:
\[
    \vec{x} = A^{-1} \cdot \vec{c}
\]
Folgende Aussagen sind äquivalent:
\begin{enumerate}
    \item $A$ ist invertierbar
    \item $A \cdot \vec{x} = \vec{c}$ hat genau eine Lösung
    \item $\operatorname{rg}(A) = n$
\end{enumerate}

\textbf{Homogenes LGS:} $A \cdot \vec{x} = \vec{0}$. Wenn $A$ invertierbar: eindeutige Lösung $\vec{x} = \vec{0}$.

\subsubsection{Gauss-Jordan-Verfahren}
Inverse berechnen durch Zeilenumformungen:
\[
    \left( A \mid E \right) \xrightarrow{\text{Gauss-Jordan}} \left( E \mid A^{-1} \right)
\]
Falls $\operatorname{rg}(A) < n$: $A$ ist nicht invertierbar.

Diagonalmatrix: $D^{-1} = \operatorname{diag}\!\left(\frac{1}{a_{11}}, \ldots, \frac{1}{a_{nn}}\right)$, falls alle $a_{ii} \neq 0$.

\textbf{Bemerkung zum Rechnen:}
\begin{itemize}
    \item Reihenfolge beibehalten: $X \cdot A + X \cdot B = X(A+B)$, aber $\neq (A+B) \cdot X$
    \item $2X + AX \neq (2+A)X$, sondern $= (2E + A)X$
\end{itemize}

\subsection{Determinanten}

\subsubsection{2x2-Matrix}
\[
    \det\begin{pmatrix} a & b \\ c & d \end{pmatrix} = ad - bc
\]

\subsubsection{3x3-Matrix (Sarrus)}
\begin{align*}
    \det(A) = &a_{11}a_{22}a_{33} + a_{12}a_{23}a_{31} + a_{13}a_{21}a_{32}\\ 
        - &a_{13}a_{22}a_{31} - a_{11}a_{23}a_{32} - a_{12}a_{21}a_{33}
\end{align*}
Sarrus gilt \textbf{nur} für $2 \times 2$ und $3 \times 3$!

\subsubsection{Laplace-Entwicklung (\(n \times n\))}
Entwicklung nach $i$-ter Zeile:
\[
    \det(A) = \sum_{j=1}^{n} (-1)^{i+j} \cdot a_{ij} \cdot \det(A_{ij})
\]
$A_{ij}$: Matrix $A$ ohne $i$-te Zeile und $j$-te Spalte. \\
\textbf{Tipp:} Nach Zeile/Spalte mit den meisten Nullen entwickeln.

\subsubsection{Geometrische Interpretation}
\begin{itemize}
    \item $|\det(A_{2\times2})|$ = Fläche des Parallelogramms der Spaltenvektoren
    \item $|\det(A_{3\times3})|$ = Volumen des Spats der Spaltenvektoren
\end{itemize}

\subsubsection{Eigenschaften der Determinante}
\begin{enumerate}
    \item $\det(E) = 1$
    \item Dreiecksmatrix: $\det(U) = u_{11} \cdot u_{22} \cdots u_{nn}$
    \item $\det(A^T) = \det(A)$
    \item $\det(A \cdot B) = \det(A) \cdot \det(B)$
    \item $\det(A^{-1}) = \frac{1}{\det(A)}$
    \item $\det(\lambda A) = \lambda^n \cdot \det(A)$ \quad ($n \times n$-Matrix)
\end{enumerate}

\subsubsection{Äquivalenzsatz}
Für eine $n \times n$-Matrix $A$ sind äquivalent:
\begin{enumerate}
    \item $\det(A) \neq 0$
    \item Spalten von $A$ sind linear unabhängig
    \item Zeilen von $A$ sind linear unabhängig
    \item $\operatorname{rg}(A) = n$
    \item $A$ ist invertierbar
    \item $A \cdot \vec{x} = \vec{c}$ hat eine eindeutige Lösung
\end{enumerate}